\documentclass[11pt,a4paper]{beamer}
\usetheme{Madrid}
\usepackage[utf8]{inputenc}
\usepackage{amsmath}
\usepackage{amsfonts}
\usepackage{amssymb}
\usepackage{graphicx}
\usepackage{xcolor}
\usepackage{tikz}
\usepackage{booktabs}
\usepackage{array}
\usepackage{colortbl}
\usepackage{multicol}

% ============================================================
% EXTREMELY SMALL FONT - MAXIMUM CONTENT PER SLIDE
% ============================================================
\usefonttheme{professionalfonts}

% Keep things compact but stable
\setbeamerfont{normal text}{size=\scriptsize}

\setbeamerfont{itemize/enumerate body}{size=\scriptsize}
\setbeamerfont{itemize/enumerate subbody}{size=\tiny}

\setbeamerfont{block body}{size=\scriptsize}
\setbeamerfont{block title}{size=\scriptsize}

\setbeamerfont{frametitle}{size=\small}
\setbeamerfont{title}{size=\normalsize}
\setbeamerfont{subtitle}{size=\small}

\setbeamerfont{author}{size=\tiny}
\setbeamerfont{date}{size=\tiny}


\renewcommand{\scriptsize}{\fontsize{6pt}{7pt}\selectfont}
\renewcommand{\footnotesize}{\fontsize{6.5pt}{8pt}\selectfont}
\renewcommand{\small}{\fontsize{7.5pt}{9pt}\selectfont}
\renewcommand{\normalsize}{\fontsize{8.5pt}{10.5pt}\selectfont}

\newcommand{\redflag}[1]{\textcolor{red}{\textbf{[!] #1}}}
\newcommand{\bluenote}[1]{\textcolor{blue}{\textbf{[i] #1}}}
\newcommand{\greennote}[1]{\textcolor{green!50!black}{\textbf{[Right] #1}}}
\newcommand{\examfocus}[1]{\textcolor{orange!80!black}{\textbf{[FOCUS] #1}}}
\newcommand{\trap}[1]{\textcolor{purple}{\textbf{[TRAP] #1}}}
\newcommand{\correct}[1]{\textcolor{green!50!black}{\textbf{[CORRECT] #1}}}
\newcommand{\scenario}[1]{\textcolor{blue!70!black}{\textbf{[SCENARIO] #1}}}
\newcommand{\keyterm}[1]{\textcolor{red!80!black}{\textbf{[KEY] #1}}}
\newcommand{\lawcite}[1]{\textcolor{brown!70!black}{\textbf{[LAW] #1}}}

\title[CAMS Domain B]{CAMS Exam Preparation\\Domain B: Global AFC Frameworks,\\Governance, and Regulations (20 percent of Exam)}
\author{Advanced AML Training Framework}
\date{}

\setbeamercovered{transparent}
\setbeamertemplate{navigation symbols}{}
\setbeamertemplate{frametitle continuation}{}
\setlength{\parskip}{2pt}
\setlength{\itemsep}{1pt}

\begin{document}
	
	% ============================================================
	% TITLE SLIDE
	% ============================================================
	\begin{frame}
		\titlepage
		\centering
		\tiny Domain B: 20 percent of CAMS Exam | 100+ Slides | FATF, UN, FIUs, Regulators, PPPs, Sanctions
	\end{frame}
	
	% ============================================================
	% SECTION 1: FATF - ROLE, FUNCTION, AND RECOMMENDATIONS
	% ============================================================
	
	\begin{frame}{B.1: FATF - Overview and Core Function}
		\begin{block}{What is FATF?}
			The Financial Action Task Force is the global standard-setter for AML/CFT.
		\end{block}
		\textbf{Key Facts:}
		\begin{itemize}
			\item Established in 1989 by the G7 at the Paris Summit
			\item Secretariat located at the OECD in Paris, France
			\item 39 members (including 2 regional organizations: European Commission and Gulf Cooperation Council)
			\item Members represent most major financial centers worldwide
			\item Sets the 40 Recommendations - global AML/CFT standards
			\item Conducts Mutual Evaluations to assess member compliance
			\item Identifies high-risk jurisdictions (Grey List and Black List)
		\end{itemize}
		\textbf{FATF's Core Functions:}
		\begin{enumerate}
			\item Set international AML/CFT standards (40 Recommendations)
			\item Assess member countries' compliance (Mutual Evaluations)
			\item Identify jurisdictions with strategic deficiencies (Grey/Black Lists)
			\item Issue guidance on emerging risks (virtual assets, DeFi, NPOs)
			\item Promote global implementation of standards
		\end{enumerate}
		\redflag{FATF Recommendations are NOT legally binding treaties, but countries must implement them to avoid economic consequences (blacklisting).}
		\examfocus{Exam tests knowledge of FATF's role, Mutual Evaluation process, and the Grey/Black Lists.}
	\end{frame}
	
	\begin{frame}{B.1.1: The 40 FATF Recommendations - Structure}
		\begin{block}{Seven Categories of the 40 Recommendations}
			The 40 Recommendations are organized into seven thematic categories.
		\end{block}
		\begin{tabular}{p{0.3\textwidth}|p{0.65\textwidth}}
			\hline
			\rowcolor{lightgray}
			\textbf{Category} & \textbf{Recommendations} \\
			\hline
			AML/CFT Policies and Coordination & Rec. 1-2 (Risk assessment, national cooperation) \\
			\hline
			Money Laundering and Confiscation & Rec. 3-4 (ML offense, confiscation) \\
			\hline
			Terrorist Financing and Proliferation & Rec. 5-8 (TF offense, sanctions, NPOs, PF) \\
			\hline
			Preventive Measures & Rec. 9-23 (CDD, recordkeeping, PEPs, DNFBPs, wire transfer) \\
			\hline
			Transparency of Legal Persons/Arrangements & Rec. 24-25 (UBO for companies and trusts) \\
			\hline
			Powers of Competent Authorities & Rec. 26-35 (supervision, FIUs, LE powers, statistics) \\
			\hline
			International Cooperation & Rec. 36-40 (MLATs, extradition, information sharing) \\
			\hline
		\end{tabular}
		\vspace{0.2cm}
		\redflag{Technical Compliance assesses whether laws exist on paper. Effectiveness assesses whether laws achieve intended outcomes.}
	\end{frame}
	
	\begin{frame}{B.1.2: FATF Mutual Evaluation Process}
		\begin{block}{How FATF Assesses Countries}
			Countries undergo peer review evaluations every 5-8 years.
		\end{block}
		\textbf{Two Pillars of Assessment:}
		\begin{itemize}
			\item Technical Compliance: Do laws exist on paper? Rated Compliant (C), Largely Compliant (LC), Partially Compliant (PC), or Non-Compliant (NC) for each of the 40 Recommendations.
			\item Effectiveness: Do laws achieve intended outcomes? Rated High, Substantial, Moderate, or Low across 11 Immediate Outcomes (IO.1 - IO.11).
		\end{itemize}
		\textbf{The 11 Immediate Outcomes (IOs):}
		\begin{enumerate}
			\item IO.1: ML/TF risks understood and coordinated domestically
			\item IO.2: International cooperation used
			\item IO.3: Supervision of FIs and DNFBPs
			\item IO.4: Preventive measures applied
			\item IO.5: Legal persons prevented from misuse
			\item IO.6: Financial intelligence used by authorities
			\item IO.7: ML investigated and prosecuted
			\item IO.8: Proceeds of crime confiscated
			\item IO.9: TF investigated and prosecuted
			\item IO.10: TF asset freezing and confiscation
			\item IO.11: Counter-proliferation financing measures
		\end{enumerate}
		\redflag{A country can have perfect Technical Compliance (laws on paper) but Low Effectiveness (no prosecutions).}
	\end{frame}
	
	\begin{frame}{B.1.3: Technical Compliance vs. Effectiveness - Exam Critical}
		\scenario{Perfect Laws, Zero Convictions}
		Country Meridian has enacted a comprehensive AML/CFT legal framework that maps directly to all 40 FATF Recommendations. The country has a functioning FIU that receives over 30,000 STRs annually. However, the FIU's financial intelligence has never been disseminated to law enforcement; the country has zero money laundering convictions in five years; and no assets have ever been confiscated. The country argues that because the laws exist on paper, it should receive a high rating.
		
		\vspace{0.1cm}
		\textbf{What the Exam Tests:} The distinction between Technical Compliance and Effectiveness.
		
		\vspace{0.1cm}
		\textbf{How to Analyze:}
		\begin{enumerate}
			\item Technical Compliance assesses whether laws exist (laws are present - likely Compliant)
			\item Effectiveness assesses whether laws achieve outcomes (no prosecutions, no confiscation - Low)
			\item High Technical Compliance does NOT compensate for Low Effectiveness
			\item Country will receive Low ratings on IO.6, IO.7, and IO.8
		\end{enumerate}
		
		\trap{Common Distractor: "A country cannot receive a low effectiveness rating if it has perfect technical compliance."}
		\correct{Correct Answer: The country will receive Compliant on Technical Compliance but Low on Effectiveness. Technical compliance and effectiveness are separate ratings.}
	\end{frame}
	
	\begin{frame}{B.1.4: FATF Grey List vs. Black List - Key Differences}
		\begin{block}{Grey List (Increased Monitoring)}
			Jurisdictions with strategic AML/CFT deficiencies that have committed to an action plan.
		\end{block}
		\textbf{Grey List Characteristics:}
		\begin{itemize}
			\item Public statement identifying jurisdictions under increased monitoring
			\item Jurisdiction has committed to resolve deficiencies within agreed timeframe
			\item FATF encourages members to apply Enhanced Due Diligence (EDD)
			\item NOT automatic termination of relationships
			\item Example jurisdictions: Turkey, Jamaica, Albania, Bulgaria (2024-2025)
		\end{itemize}
		\begin{block}{Black List (Call for Action)}
			High-risk jurisdictions subject to a call for action.
		\end{block}
		\textbf{Black List Characteristics:}
		\begin{itemize}
			\item Jurisdictions with serious strategic deficiencies
			\item FATF calls on members to apply countermeasures
			\item Countermeasures may include: terminating correspondent relationships, restricting transactions, enhanced reporting
			\item More severe consequences than Grey List
			\item Current Black List: Iran, North Korea, Myanmar (2024-2025)
		\end{itemize}
		\redflag{Key exam point: Grey List requires EDD; Black List requires countermeasures (stronger action).}
	\end{frame}
	
	\begin{frame}{B.1.5: Grey List vs. Black List - Institutional Response}
		\scenario{Bank Has Relationships in Grey List and Black List Jurisdictions}
		A multinational bank maintains correspondent accounts with banks in two jurisdictions: Country Alpha (recently placed on FATF Grey List) and Country Beta (on FATF Black List). The bank's risk committee is determining the appropriate response to each.
		
		\vspace{0.1cm}
		\textbf{What the Exam Tests:} The different obligations for Grey List vs. Black List jurisdictions.
		
		\vspace{0.1cm}
		\textbf{How to Analyze:}
		\begin{enumerate}
			\item For Country Alpha (Grey List): Apply EDD, but not automatically required to terminate
			\item For Country Beta (Black List): FATF call for countermeasures means evaluate and implement appropriate countermeasures (may include termination)
			\item Individual risk assessments required for each relationship (no blanket termination)
			\item Grey List jurisdiction has committed to resolution - may justify continued relationships with EDD
		\end{enumerate}
		
		\trap{Common Distractor: "Both jurisdictions require identical treatment - all accounts must be closed immediately."}
		\correct{Correct Answer: Grey List requires EDD; Black List requires countermeasures (stronger action). Individual risk assessments are required for both.}
	\end{frame}
	
	\begin{frame}{B.1.6: FATF Recommendation 1 - Risk-Based Approach}
		\begin{block}{Rec. 1 - Core Principle}
			Countries and financial institutions must identify, assess, and understand ML/TF risks, then apply measures commensurate with those risks.
		\end{block}
		\textbf{Key Principles:}
		\begin{itemize}
			\item Simplified measures where risks are lower (Simplified Due Diligence - SDD)
			\item Enhanced measures where risks are higher (Enhanced Due Diligence - EDD)
			\item Risk assessments must be updated regularly
			\item No one-size-fits-all approach
			\item Blanket de-risking (terminating entire categories without individual assessment) violates RBA
		\end{itemize}
		\textbf{Application Examples:}
		\begin{itemize}
			\item Low-risk customer (publicly traded company): SDD may be appropriate
			\item High-risk customer (foreign PEP): EDD required (source of wealth, source of funds, senior management approval)
		\end{itemize}
		\redflag{Not allowed: One-size-fits-all approach or blanket de-risking without individual assessment.}
	\end{frame}
	
	\begin{frame}{B.1.7: FATF Recommendation 2 - National Cooperation}
		\begin{block}{Rec. 2 - Domestic Coordination}
			Countries must establish national AML/CFT policies and ensure cooperation among competent authorities.
		\end{block}
		\textbf{Key Requirements:}
		\begin{itemize}
			\item Designate competent authorities (FIU, supervisors, law enforcement)
			\item Ensure cooperation and coordination among them
			\item Provide adequate resources and powers
			\item Establish mechanisms for information sharing
			\item Regular coordination meetings
		\end{itemize}
		\textbf{Deficiency Example (Exam Scenario):}
		A financial supervisor identifies significant AML deficiencies at several banks but never communicates these findings to the FIU or law enforcement. The supervisor issues fines but does not coordinate with other authorities. This violates Recommendation 2.
		
		\trap{Common Distractor: "Each authority operates independently; coordination is optional."}
		\correct{Correct Answer: Competent authorities must coordinate and share information. A supervisor cannot keep findings from the FIU or law enforcement.}
	\end{frame}
	
	\begin{frame}{B.1.8: FATF Recommendation 3 - Money Laundering Offense}
		\begin{block}{Rec. 3 - Criminalization of ML}
			Countries must criminalize money laundering based on the widest range of predicate offenses.
		\end{block}
		\textbf{Key Requirements:}
		\begin{itemize}
			\item Apply to all serious predicate offenses (the "all-crimes approach" is preferred)
			\item Self-laundering MUST be criminalized (predicate offender can be charged with ML)
			\item Apply regardless of where the predicate offense occurred
			\item Minimum predicate offenses include: drug trafficking, fraud, corruption, tax crimes
		\end{itemize}
		\textbf{Self-Laundering Example:}
		A drug trafficker who processes their own proceeds through shell companies can be charged with BOTH drug trafficking AND money laundering. Self-laundering is a separate criminal offense.
		
		\trap{Common Distractor: "Money laundering requires a third-party launderer separate from the predicate offender."}
		\correct{Correct Answer: FATF Recommendation 3 explicitly requires criminalization of self-laundering. The same person who committed the predicate offense can also be charged with money laundering.}
	\end{frame}
	
	\begin{frame}{B.1.9: FATF Recommendation 5 - Terrorist Financing Offense}
		\begin{block}{Rec. 5 - Criminalization of TF}
			Countries must criminalize terrorist financing as a standalone offense.
		\end{block}
		\textbf{Key Requirements:}
		\begin{itemize}
			\item Criminalize TF regardless of whether funds are actually used for a terrorist act
			\item No minimum monetary threshold (even \$1 is sufficient)
			\item Collecting or providing funds for terrorist purposes is enough
			\item Apply to nationals and entities
		\end{itemize}
		\textbf{Deficiency Example (Exam Scenario):}
		A country criminalizes terrorist financing only when the funds exceed \$10,000 and only if the funds are actually used for a specific terrorist act. This violates Recommendation 5 because:
		\begin{enumerate}
			\item No minimum threshold is permitted - any amount qualifies
			\item Collection for terrorist purposes is sufficient - actual use not required
		\end{enumerate}
		
		\trap{Common Distractor: "A \$10,000 threshold is reasonable for TF offenses."}
		\correct{Correct Answer: FATF Recommendation 5 explicitly prohibits minimum monetary thresholds for TF. Even \$1 collected for terrorist purposes is a crime.}
	\end{frame}
	
	\begin{frame}{B.1.10: FATF Recommendation 8 - Non-Profit Organizations}
		\begin{block}{Rec. 8 - NPO Vulnerability to TF}
			Countries must apply targeted risk-based supervision of NPOs vulnerable to terrorist financing abuse.
		\end{block}
		\textbf{Key Requirements:}
		\begin{itemize}
			\item Mechanisms to ensure NPOs are not misused by terrorist organizations
			\item Outreach and awareness programs for NPOs
			\item Monitoring of cross-border transactions involving NPOs
			\item Risk-based approach (not all NPOs, only vulnerable ones)
		\end{itemize}
		\textbf{Red Flags for NPOs (Exam Scenarios):}
		\begin{itemize}
			\item Large donation from a Grey List or conflict zone jurisdiction
			\item Rapid cash withdrawal (80 percent+ within 48 hours) far from headquarters
			\item Inconsistent expenditures (charitable mission vs. actual spending)
			\item No verifiable charitable activities
			\item Transfers to other NPOs in conflict zones where terrorist groups operate
		\end{itemize}
		\trap{Common Distractor: "All NPOs must be subject to the same AML controls as banks."}
		\correct{Correct Answer: Recommendation 8 requires a TARGETED RISK-BASED approach - only NPOs vulnerable to TF abuse require enhanced measures.}
	\end{frame}
	
	\begin{frame}{B.1.11: FATF Recommendation 10 - Customer Due Diligence}
		\begin{block}{Rec. 10 - CDD Requirements}
			Financial institutions must conduct CDD in specified circumstances.
		\end{block}
		\textbf{Four CDD Trigger Events:}
		\begin{enumerate}
			\item When establishing a business relationship
			\item When carrying out an occasional transaction above the threshold (typically USD/EUR 15,000)
			\item When there is a suspicion of money laundering or terrorist financing (regardless of amount)
			\item When there is doubt about previously obtained identification documents
		\end{enumerate}
		\textbf{Four Required CDD Elements:}
		\begin{enumerate}
			\item Identify the customer and verify identity using reliable, independent source documents
			\item Identify the beneficial owner and take reasonable measures to verify identity
			\item Understand the purpose and intended nature of the business relationship
			\item Conduct ongoing due diligence and scrutinize transactions throughout the relationship
		\end{enumerate}
		\redflag{All four elements are mandatory. CDD is required even for very small transactions if suspicion exists.}
	\end{frame}
	
	\begin{frame}{B.1.12: FATF Recommendation 11 - Recordkeeping}
		\begin{block}{Rec. 11 - Record Retention}
			All records must be retained for a minimum of 5 years.
		\end{block}
		\textbf{Retention Requirements:}
		\begin{itemize}
			\item Transaction records (national and international): 5 years following transaction completion
			\item Customer identification records (CDD files): 5 years after business relationship ends
			\item STRs and supporting documentation: 5 years after filing
			\item Records must be available to competent authorities upon request
		\end{itemize}
		\textbf{What Must Be Retained:}
		\begin{itemize}
			\item Copies of identification documents
			\item Account opening records
			\item Transaction history
			\item Correspondence with customer
			\item STR filings and supporting documentation
			\item Investigation notes and decisions
		\end{itemize}
		\trap{Common Distractor: "Records can be deleted when the customer closes the account."}
		\correct{Correct Answer: Minimum 5 years retention from transaction date OR account closure, whichever is later.}
	\end{frame}
	
	\begin{frame}{B.1.13: FATF Recommendation 16 - Wire Transfer Rule}
		\begin{block}{Rec. 16 - Travel Rule}
			Originator and beneficiary information must accompany wire transfers above the threshold.
		\end{block}
		\textbf{Required Information for Cross-Border Wires Above USD/EUR 1,000:}
		\begin{itemize}
			\item Originator's full name
			\item Originator's account number (or unique transaction reference)
			\item Originator's address, national ID number, customer ID number, OR date and place of birth (at least ONE)
			\item Beneficiary's full name
			\item Beneficiary's account number
		\end{itemize}
		\textbf{For Transfers Below USD/EUR 1,000:}
		\begin{itemize}
			\item Originator's name and account number required
			\item Full address/ID not required for sub-threshold transfers
		\end{itemize}
		\textbf{Receiving Institution Obligations:}
		\begin{itemize}
			\item Must obtain and hold required information
			\item Must have effective risk-based procedures to identify missing information
			\item Must take action if information not provided
		\end{itemize}
		\trap{Common Distractor: "Only the originator's name is required for wire transfers."}
		\correct{Correct Answer: Originator name, account number, and at least one additional identifier (address, ID, or date of birth) is required for transfers above USD/EUR 1,000.}
	\end{frame}
	
	\begin{frame}{B.1.14: FATF Recommendation 19 - High-Risk Jurisdictions}
		\begin{block}{Rec. 19 - Countermeasures}
			Countries must apply enhanced due diligence to high-risk jurisdictions.
		\end{block}
		\textbf{Obligations for High-Risk Jurisdictions (Black List):}
		\begin{itemize}
			\item Apply enhanced due diligence measures
			\item FATF calls on members to apply countermeasures
			\item Countermeasures may include: prohibiting correspondent relationships, restricting transactions, enhanced reporting
			\item Countermeasures are mandatory for FATF members (implemented through national laws)
		\end{itemize}
		\textbf{Obligations for Jurisdictions Under Increased Monitoring (Grey List):}
		\begin{itemize}
			\item Apply enhanced due diligence
			\item Countermeasures not automatically required (risk-based)
			\item Individual risk assessments required
			\item No blanket termination without individual assessment
		\end{itemize}
		\redflag{Grey List: EDD required. Black List: Countermeasures required (stronger action).}
	\end{frame}
	
	\begin{frame}{B.1.15: FATF Recommendation 20 - STR Filing}
		\begin{block}{Rec. 20 - Suspicious Transaction Reporting}
			STRs must be filed immediately upon suspicion, regardless of transaction value.
		\end{block}
		\textbf{Key Principles:}
		\begin{itemize}
			\item Legal standard = REASONABLE SUSPICION (not certainty, not probable cause)
			\item No minimum monetary threshold for money laundering suspicion
			\item Do NOT delay filing to gather more evidence
			\item Do NOT investigate to point of certainty before filing
			\item File first, continue investigating if needed
			\item Cannot consider commercial impact of filing
		\end{itemize}
		\textbf{What Reasonable Suspicion Means:}
		\begin{itemize}
			\item More than a hunch, less than probable cause
			\item Specific articulable facts that suggest possible money laundering
			\item Based on training, experience, and available information
			\item Does not require knowledge of underlying predicate offense
		\end{itemize}
		\trap{Common Distractor: "We need more evidence before filing a SAR."}
		\correct{Correct Answer: SAR must be filed immediately upon reasonable suspicion. Conclusive proof is NOT required.}
	\end{frame}
	
	\begin{frame}{B.1.16: FATF Recommendation 21 - Tipping-Off}
		\begin{block}{Rec. 21 - Prohibition on Disclosure}
			FIs and their employees are prohibited from disclosing that a SAR has been filed.
		\end{block}
		\textbf{Prohibited Disclosures:}
		\begin{itemize}
			\item Telling a customer a SAR has been filed on their account
			\item Disclosing that "compliance is reviewing your account" (implies SAR)
			\item Disclosing investigation to unauthorized third parties
			\item Publicly announcing a customer is under investigation
		\end{itemize}
		\textbf{Permissible Disclosures (Exceptions):}
		\begin{itemize}
			\item Internal sharing with legal department on need-to-know basis
			\item 314(b) sharing with other opted-in institutions (US specific)
			\item Responding to law enforcement subpoena
			\item Reporting to FIU (the filing itself)
		\end{itemize}
		\trap{Common Distractor: "No SAR was filed yet, so no tipping-off occurred."}
		\correct{Correct Answer: Tipping-off includes disclosing the EXISTENCE of an internal investigation, even if no SAR has been filed yet.}
	\end{frame}
	
	\begin{frame}{B.1.17: FATF Recommendation 24 - Beneficial Ownership of Legal Persons}
		\begin{block}{Rec. 24 - Corporate Transparency}
			Countries must ensure adequate, accurate, and current beneficial ownership information is obtained and held.
		\end{block}
		\textbf{Key Requirements:}
		\begin{itemize}
			\item Prohibit bearer shares OR require immobilization with a regulated custodian
			\item Beneficial ownership information must be accessible to competent authorities
			\item Sanctions for non-compliance
			\item Mechanisms to ensure information is accurate and updated
		\end{itemize}
		\textbf{Bearer Shares Vulnerability:}
		Bearer shares transfer ownership through physical delivery of the stock certificate - no public registry. This creates anonymity for the true owner. FATF requires countries to either prohibit bearer shares or require immobilization with a custodian who can identify the holder.
		
		\trap{Common Distractor: "Immobilization of bearer shares eliminates the need for UBO identification."}
		\correct{Correct Answer: The bank must verify that the custodian can identify the bearer share holder at all times AND the bank still needs to identify the beneficial owner for CDD purposes.}
	\end{frame}
	
	\begin{frame}{B.1.18: FATF Recommendation 25 - Beneficial Ownership of Trusts}
		\begin{block}{Rec. 25 - Trust Transparency}
			Trustees must hold adequate, accurate, and current beneficial ownership information on trusts.
		\end{block}
		\textbf{Who Must Be Identified for a Trust:}
		\begin{itemize}
			\item Settlor (the person who creates the trust)
			\item Trustee(s) (and UBOs if corporate trustee)
			\item Protector (if the protector has veto rights or power to replace the trustee)
			\item Beneficiaries (to the extent they are determinable)
		\end{itemize}
		\textbf{Example Trust Structure:}
		A trust has: Settlor (Individual A), Corporate Trustee, Protector (Individual B with veto power), Beneficiaries (C, D, E each with 20 percent interest).
		Under FATF standards, the ultimate beneficial owner includes: Settlor, Protector, and any beneficiary with over 25 percent interest; plus the trustee as legal owner.
		
		\trap{Common Distractor: "Only the trustee needs to be identified for a trust."}
		\correct{Correct Answer: Trustees must identify settlor, trustee, protector, and beneficiaries. A protector with veto power IS a beneficial owner.}
	\end{frame}
	
	\begin{frame}{B.1.19: FATF Recommendation 30 - Law Enforcement Powers}
		\begin{block}{Rec. 30 - Financial Investigations}
			Law enforcement must have responsibility for conducting financial investigations into ML/TF.
		\end{block}
		\textbf{Key Requirements:}
		\begin{itemize}
			\item Law enforcement must conduct parallel financial investigations in predicate crime cases
			\item Must be able to identify, trace, freeze, and seize assets
			\item Must have power to obtain financial records expeditiously (days, not months)
			\item Specialized investigators and prosecutors for financial crime
		\end{itemize}
		\textbf{Deficiency Example (Exam Scenario):}
		A country's law enforcement cannot obtain financial records directly from banks; all requests must go through the prosecutor's office, taking 6-12 months. This violates Recommendation 30 because assets can be moved before records are obtained.
		
		\trap{Common Distractor: "A prosecutor-only approval system is acceptable as long as the law exists."}
		\correct{Correct Answer: Law enforcement must have the power to obtain financial records expeditiously, without unreasonable delay. 6-12 months is not expeditious.}
	\end{frame}
	
	\begin{frame}{B.1.20: FATF Recommendation 32 - Cash Courier Controls}
		\begin{block}{Rec. 32 - Cross-Border Currency Movement}
			Countries must have declaration or disclosure systems for cross-border currency transportation.
		\end{block}
		\textbf{Required Controls:}
		\begin{itemize}
			\item Declaration OR disclosure system for currency and bearer negotiable instruments at borders
			\item Authority to stop, restrain, and question cash couriers
			\item Sanctions for false declarations or non-disclosure (seizure, fines, criminal charges)
			\item Ability to seize cash suspected of being ML/TF related
		\end{itemize}
		\textbf{Deficiency Example (Exam Scenario):}
		A country has no declaration system. Customs intercepts 20 cash couriers carrying over \$15 million in undeclared currency but imposes no penalties because no law requires declarations. This violates Recommendation 32.
		
		\trap{Common Distractor: "Cash courier controls are optional; Recommendation 32 is advisory only."}
		\correct{Correct Answer: Countries MUST have declaration or disclosure systems and MUST have authority to impose sanctions for non-disclosure.}
	\end{frame}
	
	% ============================================================
	% SECTION 2: FATF-STYLE REGIONAL BODIES (FSRBs)
	% ============================================================
	
	\begin{frame}{B.2: FATF-Style Regional Bodies (FSRBs)}
		\begin{block}{What are FSRBs?}
			Autonomous regional bodies that promote, implement, and assess FATF Recommendations within their geographic regions.
		\end{block}
		\textbf{Major FSRBs:}
		\begin{itemize}
			\item APG: Asia/Pacific Group on Money Laundering (41 members)
			\item MONEYVAL: Council of Europe committee (33 member states)
			\item GAFILAT: Financial Action Task Force of Latin America (18 members)
			\item MENAFATF: Middle East and North Africa FATF (21 members)
			\item GIABA: Inter-Governmental Action Group against Money Laundering in West Africa (17 members)
			\item ESAAMLG: Eastern and Southern Africa Anti-Money Laundering Group (18 members)
			\item EAG: Eurasian Group (9 members)
		\end{itemize}
		\textbf{Relationship to FATF:}
		\begin{itemize}
			\item FSRBs are AUTONOMOUS (not subordinate to FATF)
			\item They use the same FATF Methodology for mutual evaluations
			\item Evaluation results are reported to FATF
			\item Significant deficiencies may lead to FATF Grey/Black List referral
		\end{itemize}
		\redflag{FSRBs conduct mutual evaluations - they are not just discussion forums.}
	\end{frame}
	
	\begin{frame}{B.2.1: FSRB Mutual Evaluation Authority}
		\scenario{Country Member of FSRB but not FATF}
		Country X is a member of a FATF-Style Regional Body (MONEYVAL) but is not a FATF member. The FSRB schedules a mutual evaluation. The country's government argues that only FATF can conduct mutual evaluations, and as a non-FATF member, it is not subject to assessment.
		
		\vspace{0.1cm}
		\textbf{What the Exam Tests:} FSRB authority to evaluate member jurisdictions.
		
		\vspace{0.1cm}
		\textbf{How to Analyze:}
		\begin{enumerate}
			\item FSRBs conduct mutual evaluations using the same FATF Methodology
			\item FSRBs have authority to evaluate their members regardless of FATF membership
			\item Evaluation results are reported to FATF
			\item Significant deficiencies may lead to FATF Grey/Black List referral
		\end{enumerate}
		
		\trap{Common Distractor: "Only FATF can conduct mutual evaluations; FSRBs are just discussion forums."}
		\correct{Correct Answer: FSRBs conduct mutual evaluations of their member jurisdictions using the same methodology as FATF. The results can lead to FATF grey listing.}
	\end{frame}
	
	\begin{frame}{B.2.2: Distinguishing FATF from FSRBs}
		\begin{block}{Key Differences}
			Understanding the relationship between FATF and FSRBs is heavily tested.
		\end{block}
		\begin{tabular}{p{0.45\textwidth}|p{0.45\textwidth}}
			\hline
			\rowcolor{lightgray}
			\textbf{FATF} & \textbf{FSRBs} \\
			\hline
			Global inter-governmental body & Regional autonomous bodies \\
			Sets the 40 Recommendations & Promote and implement Recommendations \\
			39 members & Varies by region (9-41 members) \\
			Conducts mutual evaluations of members & Conducts mutual evaluations of members \\
			Issues Grey/Black Lists & Refer findings to FATF for listing \\
			Based in Paris (OECD) & Various locations (Bangkok, Buenos Aires, etc.) \\
			\hline
		\end{tabular}
		\vspace{0.2cm}
		\textbf{Exam Point:} FSRB mutual evaluations can lead to FATF grey listing. A country does NOT need to be a FATF member to be evaluated by an FSRB or to be placed on the FATF Grey/Black Lists.
		
		\trap{Common Distractor: "FSRBs are subordinate to FATF and staffed by FATF employees."}
		\correct{Correct Answer: FSRBs are autonomous regional bodies, not subordinate to FATF. They promote and assess FATF Recommendations in their regions.}
	\end{frame}
	
	% ============================================================
	% SECTION 3: UNITED NATIONS - ROLE AND SANCTIONS
	% ============================================================
	
	\begin{frame}{B.3: United Nations - Role in AML/CFT}
		\begin{block}{UN Sanctions Framework}
			UN Security Council resolutions are binding on all UN member states.
		\end{block}
		\textbf{Key UN Sanctions Resolutions:}
		\begin{itemize}
			\item UNSCR 1267 (1999): Al-Qaida and Taliban sanctions
			\item UNSCR 1373 (2001): Counter-terrorism and terrorist financing (general)
			\item UNSCR 1540 (2004): WMD proliferation financing
			\item UNSCR 1718 (2006): North Korea (nuclear and missile programs)
			\item UNSCR 1988 (2011): Taliban sanctions (separate from Al-Qaida)
			\item UNSCR 2231 (2015): Iran nuclear program (JCPOA)
		\end{itemize}
		\textbf{Bank Obligations Under UN Sanctions:}
		\begin{itemize}
			\item Freeze assets immediately without prior notice
			\item Report freezing action to UNSC committee and FIU
			\item Prohibit funds from being made available to designated persons
			\item Screen against UN Consolidated Sanctions List
		\end{itemize}
		\redflag{UN sanctions are BINDING on all UN member states - no option to opt out.}
	\end{frame}
	
	\begin{frame}{B.3.1: UNSCR 1373 - Terrorist Financing}
		\begin{block}{UNSCR 1373 - Counter-Terrorism}
			Binding resolution requiring all UN member states to criminalize terrorist financing and freeze terrorist assets.
		\end{block}
		\textbf{Key Requirements:}
		\begin{itemize}
			\item Criminalize provision or collection of funds for terrorist purposes
			\item Freeze assets of persons who commit or support terrorist acts
			\item Prohibit funds from being made available to terrorists
			\item Prevent movement of terrorists across borders
			\item Cooperate in criminal investigations
		\end{itemize}
		\textbf{Bank Obligations:}
		\begin{itemize}
			\item Screen against UNSCR 1373 designated persons list
			\item Freeze assets immediately without prior notice
			\item Report freezing actions to competent authorities
			\item Cannot notify the designated person (tipping-off)
		\end{itemize}
		\trap{Common Distractor: "UNSCR 1373 only applies to countries, not to private financial institutions."}
		\correct{Correct Answer: UNSCR 1373 is implemented through national laws that apply directly to financial institutions. Banks must freeze assets of designated persons.}
	\end{frame}
	
	\begin{frame}{B.3.2: UNSCR 1540 - Proliferation Financing}
		\begin{block}{UNSCR 1540 - WMD Proliferation}
			Binding resolution requiring member states to prevent WMD proliferation financing.
		\end{block}
		\textbf{Key Requirements:}
		\begin{itemize}
			\item Establish effective controls over chemical, biological, radiological, and nuclear (CBRN) materials
			\item Prohibit non-state actors from manufacturing or acquiring WMD-related materials
			\item Establish controls over dual-use goods exports
			\item Criminalize proliferation financing
		\end{itemize}
		\textbf{Bank Obligations Under UNSCR 1540:}
		\begin{itemize}
			\item Screen transactions for dual-use goods exports to sanctioned jurisdictions
			\item Identify proliferation financing red flags (inconsistent end-use, routing through intermediaries)
			\item File STRs for suspicious proliferation-related transactions
			\item Enhanced due diligence for proliferation-sensitive transactions
		\end{itemize}
		\textbf{Red Flags for Proliferation Financing:}
		\begin{itemize}
			\item Dual-use goods (radiation-hardened chips, maraging steel, vacuum pumps) to sanctioned countries
			\item End-use claims inconsistent with goods (e.g., "atmospheric research" for missile components)
			\item Payment routing through multiple third-party intermediaries
			\item Unusual shipping routes or free trade zone transshipment
		\end{itemize}
		\redflag{Proliferation financing is a separate FATF offense category (Recommendation 7).}
	\end{frame}
	
	% ============================================================
	% SECTION 4: REGULATORS, LAW ENFORCEMENT, AND FIUs
	% ============================================================
	
	\begin{frame}{B.4: Financial Intelligence Units (FIUs)}
		\begin{block}{What is an FIU?}
			The national center for receiving, analyzing, and disseminating financial intelligence.
		\end{block}
		\textbf{Core FIU Functions:}
		\begin{itemize}
			\item Receive STRs from financial institutions and DNFBPs
			\item Analyze STRs to identify patterns and develop intelligence
			\item Disseminate financial intelligence to competent authorities (law enforcement, prosecutors)
			\item Exchange information with foreign FIUs (via Egmont Group)
			\item Provide feedback to reporting institutions on STR quality
			\item Conduct strategic analysis on ML/TF trends
		\end{itemize}
		\textbf{Examples of FIUs by Country:}
		\begin{itemize}
			\item United States: FinCEN (Financial Crimes Enforcement Network)
			\item United Kingdom: UKFIU (part of National Crime Agency)
			\item France: TRACFIN
			\item Germany: FIU Germany (Zentralstelle für Finanztransparenzuntersuchungen)
			\item Spain: SEPBLAC
			\item Australia: AUSTRAC
			\item Singapore: STRO (Suspicious Transaction Reporting Office)
		\end{itemize}
		\redflag{FIUs have analytical and intelligence functions, not law enforcement powers.}
	\end{frame}
	
	\begin{frame}{B.4.1: FIU Information Sharing - Egmont Group}
		\begin{block}{Egmont Group Functions}
			The Egmont Group is the global network of FIUs, not just a definition.
		\end{block}
		\textbf{Egmont Group Core Functions:}
		\begin{itemize}
			\item Secure communication platform (Egmont Secure Web) for FIU-to-FIU information exchange
			\item Established "Egmont Group Principles for Information Exchange Between FIUs"
			\item Provides training and technical assistance for member FIUs
			\item Membership requires meeting criteria (member of FATF or FSRB, operational FIU)
		\end{itemize}
		\textbf{Limitations on Egmont-Exchanged Intelligence:}
		\begin{itemize}
			\item Used only for AML/CFT purposes (cannot be used for tax investigations or commercial purposes)
			\item Receiving FIU must maintain confidentiality
			\item Cannot share with third parties without originating FIU's prior authorization
			\item NOT automatically admissible in criminal court (unlike MLAT-obtained evidence)
		\end{itemize}
		\trap{Common Distractor: "Egmont Group information exchange can replace MLATs for criminal prosecution."}
		\correct{Correct Answer: Egmont exchanges are for INTELLIGENCE sharing (not admissible in court). MLATs are required for evidence admissible in criminal proceedings.}
	\end{frame}
	
	\begin{frame}{B.4.2: Law Enforcement Cooperation - MLATs}
		\begin{block}{Mutual Legal Assistance Treaties (MLATs)}
			Formal, legally binding mechanisms for governments to exchange evidence across borders.
		\end{block}
		\textbf{MLAT Functions:}
		\begin{itemize}
			\item Exchange of evidence, judicial documents, and witness testimony
			\item Cross-border asset freezing and forfeiture
			\item Service of legal documents
			\item Execution of search warrants
		\end{itemize}
		\textbf{MLAT Process:}
		\begin{enumerate}
			\item Requesting country submits formal request to central authority
			\item Requesting country must provide evidence of criminal investigation
			\item Requested country executes request according to its laws
			\item Evidence returned to requesting country
			\item Evidence is admissible in criminal proceedings
		\end{enumerate}
		\textbf{MLAT vs. Egmont Comparison:}
		\begin{itemize}
			\item MLAT: Evidence admissible in court (formal, slow, requires treaty)
			\item Egmont: Intelligence sharing (informal, fast, not court-admissible without conversion)
		\end{itemize}
		\redflag{For criminal prosecution, use MLAT. For intelligence sharing, use Egmont.}
	\end{frame}
	
	\begin{frame}{B.4.3: MLAT Process - Exam Scenario}
		\scenario{Country A Needs Bank Records from Country B for Criminal Prosecution}
		Country A has identified \$30 million in bribery proceeds deposited into a bank account in Country B. Country A's prosecutor needs bank records to support a criminal prosecution. Country A and Country B have an active MLAT and are both members of the Egmont Group.
		
		\vspace{0.1cm}
		\textbf{What the Exam Tests:} Proper channel for obtaining admissible evidence vs. intelligence.
		
		\vspace{0.1cm}
		\textbf{How to Analyze:}
		\begin{enumerate}
			\item The FIU-to-FIU channel (Egmont) provides intelligence, NOT admissible evidence
			\item For criminal prosecution, an MLAT request is required
			\item The proper approach: FIU exchanges intelligence (Egmont) PLUS law enforcement submits MLAT request for admissible evidence
			\item Cannot rely solely on Egmont for criminal prosecution
		\end{enumerate}
		
		\trap{Common Distractor: "Country A's FIU can obtain records through Egmont and use them directly in criminal court."}
		\correct{Correct Answer: MLAT is the proper channel for obtaining bank records admissible in criminal proceedings. Egmont provides intelligence only, not court-admissible evidence.}
	\end{frame}
	
	% ============================================================
	% SECTION 5: PUBLIC-PRIVATE PARTNERSHIPS (PPPs)
	% ============================================================
	
	\begin{frame}{B.5: Public-Private Partnerships (PPPs)}
		\begin{block}{What are AML/CFT PPPs?}
			Structured collaboration between government authorities (FIUs, law enforcement, regulators) and financial institutions to combat financial crime.
		\end{block}
		\textbf{PPP Objectives:}
		\begin{itemize}
			\item Improve intelligence sharing and detection of complex financial crime networks
			\item Help financial institutions understand emerging typologies
			\item Help law enforcement prioritize financial intelligence
			\item Improve SAR quality through feedback
			\item Identify systemic vulnerabilities
		\end{itemize}
		\textbf{What PPPs DO NOT Replace:}
		\begin{itemize}
			\item SAR filing obligations (banks still must file SARs)
			\item Independent testing of AML programs
			\item CDD and KYC obligations
			\item Regulatory supervision
		\end{itemize}
		\textbf{Examples of PPPs:}
		\begin{itemize}
			\item JMLIT (UK): Joint Money Laundering Intelligence Task Force
			\item FinCEN Exchange (US): Information sharing between FinCEN and financial institutions
			\item FIU-Netherlands PPP program
			\item Singapore's COSMIC platform
		\end{itemize}
		\redflag{PPPs enhance, not replace, existing AML obligations.}
	\end{frame}
	
	\begin{frame}{B.5.1: PPP Strengths and Limitations}
		\begin{block}{What PPPs Achieve}
			Understanding the value and limitations of PPPs is tested on the exam.
		\end{block}
		\textbf{Strengths of PPPs:}
		\begin{itemize}
			\item Real-time threat intelligence sharing (vs. delayed SAR processing)
			\item Identification of emerging typologies (e.g., COVID fraud, ransomware trends)
			\item Feedback loop improves SAR quality
			\item Builds trust between public and private sectors
			\item More efficient targeting of high-risk areas
		\end{itemize}
		\textbf{Limitations of PPPs:}
		\begin{itemize}
			\item Cannot replace SAR filing obligations
			\item Information sharing subject to confidentiality and privacy laws
			\item Not all institutions can participate (resource constraints)
			\item May not be available in all jurisdictions
			\item Does not eliminate need for individual institution risk assessments
		\end{itemize}
		\textbf{Data Sharing Considerations:}
		\begin{itemize}
			\item Must comply with data privacy laws (GDPR, CCPA, etc.)
			\item Need information sharing agreements (MOUs)
			\item Data minimization principles apply
			\item Confidentiality obligations must be maintained
		\end{itemize}
		\trap{Common Distractor: "PPPs replace SAR filing obligations for participating banks."}
		\correct{Correct Answer: PPPs enhance intelligence sharing but do NOT replace SAR filing obligations. Banks must still file SARs.}
	\end{frame}
	
	% ============================================================
	% SECTION 6: US AND EU REGULATIONS
	% ============================================================
	
	\begin{frame}{B.6: US AML/CFT Regulations}
		\begin{block}{US Regulatory Framework}
			The United States has a comprehensive AML/CFT regulatory structure.
		\end{block}
		\textbf{Primary US AML Laws:}
		\begin{itemize}
			\item Bank Secrecy Act (BSA) of 1970 - foundation of US AML
			\item USA PATRIOT Act of 2001 - expanded AML requirements
			\item Anti-Money Laundering Act (AMLA) of 2020 - most recent major update
			\item Corporate Transparency Act (CTA) - beneficial ownership reporting
		\end{itemize}
		\textbf{Key US Regulators:}
		\begin{itemize}
			\item FinCEN: Issues regulations, administers BSA, operates as FIU
			\item OFAC: Administers sanctions programs (US Treasury)
			\item Federal banking regulators: OCC, Federal Reserve, FDIC, NCUA
			\item SEC: Securities industry AML oversight
			\item CFTC: Commodities and derivatives oversight
			\item FinCEN 314(a) and 314(b): Information sharing programs
		\end{itemize}
		\textbf{Key USA PATRIOT Act Sections Tested:}
		\begin{itemize}
			\item Section 311: Designation of primary money laundering concern
			\item Section 313: Prohibition on correspondent accounts for shell banks
			\item Section 314(a): Law enforcement information requests
			\item Section 314(b): Voluntary information sharing among FIs
			\item Section 319(b): Record production and account seizure
		\end{itemize}
		\redflag{USA PATRIOT Act sections 311, 313, 314(a), 314(b), and 319(b) are heavily tested.}
	\end{frame}
	
	\begin{frame}{B.6.1: USA PATRIOT Act Section 311}
		\begin{block}{Section 311 - Primary Money Laundering Concern}
			Treasury Secretary can designate foreign financial institutions as primary ML concern.
		\end{block}
		\textbf{Measures That Can Be Imposed:}
		\begin{itemize}
			\item Prohibit U.S. FIs from opening or maintaining correspondent or payable-through accounts
			\item Impose special recordkeeping and reporting requirements on transactions
			\item Impose conditions on correspondent accounts (e.g., requiring pre-approval)
		\end{itemize}
		\textbf{What Section 311 Does NOT Authorize:}
		\begin{itemize}
			\item Seizing personal assets of employees
			\item Criminal indictment of all officers
			\item Automatic freezing of all customer accounts
		\end{itemize}
		\trap{Common Distractor: "Section 311 allows seizure of all assets of the designated institution."}
		\correct{Correct Answer: Section 311 allows prohibiting correspondent accounts, imposing special recordkeeping/reporting, and imposing conditions on accounts - not asset seizure.}
	\end{frame}
	
	\begin{frame}{B.6.2: USA PATRIOT Act Section 313 - Shell Bank Prohibition}
		\begin{block}{Section 313 - Definition and Prohibition}
			A shell bank is a bank with no physical presence in any country and not affiliated with a regulated financial institution.
		\end{block}
		\textbf{Absolute Prohibition:}
		U.S. banks are absolutely prohibited from maintaining correspondent accounts for foreign shell banks. NO exceptions.
		
		\textbf{Not a Shell Bank (Acceptable):}
		\begin{itemize}
			\item Bank with physical office and staff in its home country
			\item Bank affiliated with a regulated depository institution
		\end{itemize}
		\textbf{Exam Scenario:}
		A U.S. bank discovers its foreign respondent bank has no physical office, no staff, and operates through a shared services agreement. The bank is not affiliated with any regulated institution. This is a shell bank. The U.S. bank must terminate the correspondent account immediately.
		
		\trap{Common Distractor: "The bank can continue the relationship if it files SARs every 90 days."}
		\correct{Correct Answer: Section 313 absolutely prohibits correspondent accounts for foreign shell banks. No exceptions. The account must be terminated.}
	\end{frame}
	
	\begin{frame}{B.6.3: USA PATRIOT Act Section 314(a) and 314(b)}
		\begin{block}{Section 314(a) - Mandatory Law Enforcement Requests}
			FinCEN sends requests identifying subjects of investigation; banks must search records.
		\end{block}
		\textbf{314(a) Bank Obligations:}
		\begin{itemize}
			\item Search accounts and transactions according to FinCEN instructions
			\item Search for specified time period (e.g., last 6 months)
			\item Maintain confidentiality - cannot disclose request to customer
			\item Report matches to law enforcement via FinCEN
			\item NOT required to freeze accounts automatically
		\end{itemize}
		\begin{block}{Section 314(b) - Voluntary Information Sharing}
			Financial institutions that have filed opt-in notification may voluntarily share information.
		\end{block}
		\textbf{314(b) Key Features:}
		\begin{itemize}
			\item Share information about persons suspected of ML/TF activity
			\item Receive safe harbor from liability for sharing
			\item Build complete picture across institutions
			\item May share SAR information (exception to tipping-off)
		\end{itemize}
		\redflag{314(b) sharing is explicitly permitted by law and is NOT tipping-off.}
	\end{frame}
	
	\begin{frame}{B.6.4: USA PATRIOT Act Section 319(a) and 319(b)}
		\begin{block}{Section 319(a) - Forfeiture of Funds}
			Allows U.S. government to seize funds from correspondent accounts traceable to money laundering.
		\end{block}
		\textbf{319(a) Key Points:}
		\begin{itemize}
			\item Applies to funds in correspondent accounts
			\item Funds must be traceable to money laundering
			\item Subject to judicial process
		\end{itemize}
		\begin{block}{Section 319(b) - Record Production}
			Allows U.S. government to demand that foreign banks produce records of U.S. correspondent accounts.
		\end{block}
		\textbf{319(b) Key Points:}
		\begin{itemize}
			\item Foreign bank must maintain records of its U.S. correspondent account
			\item If foreign bank refuses subpoena, U.S. bank may be required to terminate the account
			\item Funds in the account may be seized
		\end{itemize}
		\textbf{Distinction:}
		\begin{itemize}
			\item 319(a): Forfeiture of funds (asset recovery)
			\item 319(b): Record production (information gathering)
		\end{itemize}
		\trap{Common Distractor: "Section 319(a) and 319(b) apply only to FATF Black List jurisdictions."}
		\correct{Correct Answer: Sections 319(a) and 319(b) apply to all foreign banks with U.S. correspondent accounts, regardless of FATF status.}
	\end{frame}
	
	\begin{frame}{B.6.5: EU AML Directives - 4AMLD, 5AMLD, 6AMLD}
		\begin{block}{EU Fourth AML Directive (4AMLD) - 2015}
			Key provisions introduced.
		\end{block}
		\textbf{4AMLD Key Provisions:}
		\begin{itemize}
			\item Beneficial ownership registers for corporate entities
			\item Risk-based approach to AML/CFT supervision
			\item Enhanced due diligence for PEPs (both domestic and foreign)
			\item Extended AML requirements to casinos
		\end{itemize}
		\begin{block}{EU Fifth AML Directive (5AMLD) - 2018}
			Expanded and strengthened requirements.
		\end{block}
		\textbf{5AMLD Key Provisions:}
		\begin{itemize}
			\item Extended CDD to VASPs (virtual asset service providers)
			\item Publicly accessible central registries of beneficial owners
			\item Lowered prepaid card CDD threshold from EUR 250 to EUR 150
			\item Bank account registers for FIUs
		\end{itemize}
		\begin{block}{EU Sixth AML Directive (6AMLD) - 2020}
			Harmonized criminal law provisions.
		\end{block}
		\textbf{6AMLD Key Provisions:}
		\begin{itemize}
			\item Harmonized 22 predicate offenses across all EU member states
			\item Extended criminal liability to legal persons (corporations)
			\item Minimum maximum sentences: 4 years for ML, up to 15 years for aggravating circumstances
		\end{itemize}
		\redflag{6AMLD added corporate criminal liability - corporations can be prosecuted for ML.}
	\end{frame}
	
	% ============================================================
	% SECTION 7: OTHER LAWS IMPACTING AML (GDPR, ESG, ETC.)
	% ============================================================
	
	\begin{frame}{B.7: Other Laws Impacting AML Applications}
		\begin{block}{Intersection of AML with Other Regulatory Areas}
			AML compliance does not operate in isolation; other laws create overlapping obligations.
		\end{block}
		\textbf{Key Overlapping Regulatory Areas:}
		\begin{itemize}
			\item Data Privacy (GDPR, CCPA): Right to erasure vs. AML retention
			\item Consumer Protection: Unfair treatment of customers, account closures
			\item ESG (Environmental, Social, Governance): Corruption risk, human trafficking detection
			\item Conduct and Ethics: Whistleblower protections, ethical culture
			\item Financial Inclusion: De-risking impacts on underserved populations
		\end{itemize}
		\textbf{Resolving Conflicts - Exam Approach:}
		\begin{itemize}
			\item AML obligations generally override privacy rights for required data retention
			\item Data protection laws include exceptions for law enforcement and AML
			\item Financial inclusion is a consideration, not an override of AML obligations
		\end{itemize}
		\redflag{AML retention requirements override GDPR right to erasure for required records.}
	\end{frame}
	
	\begin{frame}{B.7.1: GDPR vs. AML - Resolving Conflicts}
		\scenario{Bank Receives GDPR Right to Erasure Request for SAR Customer}
		A European bank receives a GDPR "right to erasure" request from a customer demanding deletion of all personal data, including KYC documents and transaction records. The customer was the subject of a SAR filed 18 months ago. The DPO argues GDPR requires deletion. The AML Officer argues AML recordkeeping requires retention (minimum 5 years).
		
		\vspace{0.1cm}
		\textbf{What the Exam Tests:} GDPR's exception for legal obligations (AML retention).
		
		\vspace{0.1cm}
		\textbf{How to Analyze:}
		\begin{enumerate}
			\item GDPR Article 17(3) exempts data retention required by law
			\item AML recordkeeping (minimum 5 years) is a legal obligation
			\item The bank can lawfully refuse to delete AML-required records
			\item SAR confidentiality also prohibits acknowledging the SAR to the customer
		\end{enumerate}
		
		\trap{Common Distractor: "GDPR's right to erasure is absolute and overrides all other laws."}
		\correct{Correct Answer: AML retention requirements are a legal exception to GDPR erasure rights. The bank may refuse deletion for records it is required to retain.}
	\end{frame}
	
	% ============================================================
	% SECTION 8: CONCLUSION AND EXAM REMINDERS
	% ============================================================
	
	\begin{frame}{Domain B Summary: Top 30 Most Tested Concepts}
		\begin{multicols}{2}
			\begin{enumerate}
				\item FATF sets global AML standards (40 Recommendations)
				\item Mutual Evaluation assesses Technical Compliance AND Effectiveness
				\item High Technical Compliance does not guarantee high Effectiveness
				\item Grey List: EDD required; Black List: Countermeasures required
				\item FSRBs are autonomous regional bodies, not subordinate to FATF
				\item FSRBs conduct mutual evaluations using FATF methodology
				\item UN sanctions are binding on all UN member states
				\item UNSCR 1373: Terrorist financing
				\item UNSCR 1540: WMD proliferation financing
				\item FIUs receive, analyze, and disseminate STR intelligence
				\item Egmont Group: FIU-to-FIU intelligence sharing (not court-admissible)
				\item MLAT: Court-admissible evidence (formal, slow)
				\item PPPs enhance intelligence sharing, do NOT replace SAR filing
				\item USA PATRIOT Sections 311, 313, 314(a), 314(b), 319(b) heavily tested
				\item Section 313: Shell banks prohibited (no physical presence)
				\item Section 314(a): Mandatory law enforcement searches
				\item Section 314(b): Voluntary information sharing (safe harbor)
				\item 6AMLD: Corporate criminal liability for ML
				\item 5AMLD: VASPs regulated, prepaid card threshold lowered
				\item 4AMLD: BO registers, risk-based approach
				\item GDPR Article 17(3): AML retention exception
				\item STR filing standard: Reasonable suspicion (not proof)
				\item Tipping-off includes disclosing investigation exists
				\item NPO red flags: rapid cash withdrawal, conflict zone operations
				\item Proliferation financing red flags: dual-use goods, inconsistent end-use
				\item Bearer shares require immobilization AND UBO identification
				\item Trust UBOs: settlor, trustee, protector, determinable beneficiaries
				\item Law enforcement must access records expeditiously (not 6-12 months)
				\item Cash courier declaration systems required
				\item Self-laundering IS criminalized under FATF Rec. 3
			\end{enumerate}
		\end{multicols}
	\end{frame}
	
	\begin{frame}{Critical Red Flags - Domain B Quick Reference}
		\begin{multicols}{2}
			\begin{itemize}
				\item Perfect laws with zero convictions → Low Effectiveness
				\item Blanket termination of all African banks → De-risking (violates RBA)
				\item FIU receives STRs but never disseminates → Low IO.6 rating
				\item No ML convictions in 5 years → Low IO.7 rating
				\item No assets confiscated despite laws → Low IO.8 rating
				\item \$10,000 threshold for TF offense → Violates Rec. 5
				\item No BO registry or inaccessible registry → Violates Rec. 24
				\item 6-12 month delay for bank records → Violates Rec. 30
				\item No cash declaration system → Violates Rec. 32
				\item U.S. bank has account for shell bank → Violates Section 313
				\item Telling customer "compliance is reviewing" → Tipping-off
				\item Accepting missing Travel Rule data → VASP violation
				\item No EDD for Grey List jurisdiction → Rec. 19 violation
				\item Self-laundering not criminalized → Violates Rec. 3
				\item No independent audit of NPOs with conflict zone ops → Rec. 8 deficiency
				\item FIU intelligence not used by law enforcement → Low IO.6
				\item Supervisors never penalize banks → Rec. 26 deficiency
			\end{itemize}
		\end{multicols}
	\end{frame}
\begin{frame}{Exam Day Reminders for Domain B}
	\begin{itemize}
		\item Domain B is 20 percent of exam
		\item FATF is the single most heavily tested topic
		\item Know the difference between Grey List (EDD) and Black List (countermeasures)
		\item Technical Compliance vs. Effectiveness is a key distinction
		\item FSRBs conduct mutual evaluations - their findings can lead to FATF listing
		\item USA PATRIOT Act sections 311, 313, 314(a), 314(b), 319(b) are essential
		\item UN sanctions are binding - freeze assets immediately without notice
		\item Egmont = intelligence (not admissible); MLAT = evidence (admissible)
		\item PPPs enhance, not replace, SAR filing obligations
		\item 6AMLD added corporate criminal liability
		\item 5AMLD added VASPs to AML requirements
		\item GDPR has AML retention exception (Article 17(3))
		\item STR filing standard = reasonable suspicion (not proof)
		\item Tipping-off includes disclosing investigation exists
		\item Self-laundering IS criminalized under FATF standards
		\item Bearer shares require both immobilization AND UBO identification
		\item Trust UBOs include protector with veto power
	\end{itemize}
	\vspace{0.3cm}
	\centering
	\greennote{Good luck on your CAMS exam!}
\end{frame}

\end{document}